\begin{table}[tb]
\centering
\caption{Annotation guidelines for checklist items and grader assignment.}
\label{tab:annotation_checklists}
\small
\begin{tabular}{p{0.18\linewidth}p{0.37\linewidth}p{0.35\linewidth}}
\toprule[1.25pt]
\textbf{Aspect} & \textbf{Do} $\checkmark$ & \textbf{Avoid} $\times$ \\
\midrule[0.75pt]
\textsc{Verifiability} &
Each item should be checkable from the trajectory, produced artifacts, or tool records. &
Do not include criteria that require unstated assumptions or external judgment beyond the available evidence. \\
\midrule[0.75pt]
\textsc{Atomicity} &
Each item should cover one requirement whenever possible. &
Do not combine several independent requirements into one checklist item. \\
\midrule[0.75pt]
\textsc{Outcome focus} &
Write checklist items as final completion conditions. &
Do not use checklist items to score whether the agent was proactive during the interaction. \\
\midrule[0.75pt]
\textsc{Grader choice} &
Use rubric-based evaluation for semantic criteria and rule-based verification for exact structured criteria. &
Do not use an LLM rubric when an exact programmatic check is available and reliable. \\
\midrule[0.75pt]
\textsc{Auditability} &
Make the expected evidence clear enough for another expert to inspect the judgment. &
Do not write criteria whose satisfaction cannot be traced to concrete evidence. \\
\bottomrule[1.25pt]
\end{tabular}
\end{table}